\documentclass[a4paper]{article}
% Español (México) con XeLaTeX: fontspec para fuentes Unicode, polyglossia
% para la división silábica y los nombres del documento en la variante mexicana.
\usepackage{fontspec}
\usepackage{polyglossia}
\setdefaultlanguage[variant=mexican]{spanish}
% Los nombres chinos van como «pinyin (original)»: xeCJK da a los caracteres
% Han su propia fuente; sin ella XeLaTeX los omite con un aviso.
% FandolSong no cubre algunos caracteres de nombres (U+73FA); WenQuanYi Zen Hei sí.
\usepackage[AutoFallBack=true]{xeCJK}
\setCJKmainfont{FandolSong-Regular.otf}
\setCJKfallbackfamilyfont{\CJKrmdefault}{WenQuanYi Zen Hei}
% polyglossia no traduce los nombres de listings.
\AtBeginDocument{\providecommand{\lstlistingname}{}\renewcommand{\lstlistingname}{Listado}%
  \providecommand{\lstlistlistingname}{}\renewcommand{\lstlistlistingname}{Índice de listados}}
\usepackage{amsmath, amssymb}
\usepackage{graphicx}
\usepackage[margin=2.35cm]{geometry}
\usepackage[most]{tcolorbox}
\usepackage{booktabs}
\usepackage{tabularx}
\usepackage{array}
\usepackage{float}
\usepackage{xcolor}
\usepackage{hyperref}
\usepackage{fancyhdr}
\setCJKmainfont[AutoFakeBold=2.4]{AR PL UMing CN}
\setCJKsansfont[AutoFakeBold=2.4]{Droid Sans Fallback}
\setCJKmonofont[AutoFakeBold=2.4]{Droid Sans Fallback}
\hypersetup{colorlinks=true, linkcolor=blue!55!black, urlcolor=blue!60!black}
\graphicspath{{./}{figures/}}
\setlength{\parskip}{0.55em}
\setlength{\parindent}{2em}
\setlength{\emergencystretch}{3em}
\renewcommand{\arraystretch}{1.18}
\newtcolorbox{knowledgebox}[1]{enhanced, colback=blue!5!white, colframe=blue!70!black, colbacktitle=blue!70!black, coltitle=white, fonttitle=\bfseries, title={#1}, sharp corners, breakable}
\newtcolorbox{importantbox}[1]{enhanced, colback=yellow!10!white, colframe=yellow!75!black, colbacktitle=yellow!75!black, coltitle=black, fonttitle=\bfseries, title={#1}, sharp corners, breakable}
\newtcolorbox{warningbox}[1]{enhanced, colback=red!5!white, colframe=red!70!black, colbacktitle=red!70!black, coltitle=white, fonttitle=\bfseries, title={#1}, sharp corners, breakable}
\pagestyle{fancy}
\fancyhf{}
\fancyfoot[C]{\thepage}
\fancyfoot[R]{\footnotesize\color{gray}\href{https://github.com/hqhq1025/ai-course-notes}{AI Course Notes}}
\renewcommand{\headrulewidth}{0pt}
\renewcommand{\footrulewidth}{0pt}
\newcommand{\videourl}{https://www.youtube.com/watch?v=vWrYHvSRz0s}
\newcommand{\repourl}{https://github.com/hqhq1025/ai-course-notes}

\begin{document}
\begin{titlepage}
\centering
{\Large Notas didácticas de una entrevista a fondo\par}
\vspace{1.0cm}
{\huge\bfseries Diez años de investigación multimodal, la inspiración de o1 y el próximo momento GPT-4}\par
\vspace{0.5cm}
{\Large Zhang Xiaojun Podcast: Li Guangmi (李广密), co-host, conversa con Zhang Xiangyu (张祥雨)\par}
\vspace{0.35cm}
{\large Elaboradas a partir del video público, la estructura de capítulos, la transcripción, la descripción del video y diagramas conceptuales propios\par}
\vspace{0.3cm}
{\large 2025-06-03\par}
\vspace{0.85cm}
\includegraphics[width=0.88\textwidth,height=0.42\textheight,keepaspectratio]{cover.jpg}\par
\vfill
\begin{tcolorbox}[width=0.92\textwidth, colback=black!2!white, colframe=black!55, sharp corners]
\textbf{Título del video}: 102. Conversación con Zhang Xiangyu: las dificultades de la investigación multimodal y los dos «momentos GPT-4» de los próximos dos años\par
\textbf{Canal del video}: Zhang Xiaojun Podcast\par
\textbf{Duración del video}: 02:28:44\par
\textbf{Enlace del video}: \href{\videourl}{\nolinkurl{\videourl}}\par
\textbf{Nota sobre los materiales}: YouTube no ofrece subtítulos descargables; el texto principal se elaboró a partir de una transcripción por bloques con faster-whisper large-v3 (CPU, int8), la descripción del video, la estructura de capítulos y figuras didácticas propias. Las tomas fijas de la entrevista no se repiten en el texto principal.\par
\textbf{Página del proyecto}: \href{\repourl}{\nolinkurl{\repourl}}\par
\end{tcolorbox}
\end{titlepage}

\begingroup
\small
\setlength{\parskip}{0pt}
\renewcommand{\baselinestretch}{0.92}\selectfont
\tableofcontents
\endgroup
\newpage

\section{Introducción: por qué este episodio es una clase sobre la trayectoria técnica multimodal}

Esta sección establece primero un marco de lectura. Este episodio es una entrevista de alta densidad técnica, poco común en el programa de entrevistas de negocios de Zhang Xiaojun (张小珺). Li Guangmi (李广密), como co-host, centra las preguntas en diez años de historia multimodal, las diferencias entre visión y lenguaje, lo que el paradigma de o1 sugiere para lo multimodal y el momento GPT-4 que podría llegar en los próximos dos años. El relato de Zhang Xiangyu (张祥雨) no es una colección de anécdotas sobre artículos, sino una trayectoria técnica: la CV aprendió de manera continua del NLP, pero la imagen estática en sí es distinta del lenguaje natural; el verdadero avance multimodal quizá requiera volver a entender next token prediction, RL, CoT, el espacio de acciones y el aprendizaje en línea.

\begin{importantbox}{Tesis central del episodio}
La dificultad de lo multimodal no se reduce a «poner imágenes, video y texto en un mismo modelo». En la imagen estática, la generación, la comprensión y la alineación con los humanos están separadas por naturaleza; la generación visual y la comprensión visual llevan mucho tiempo sin fusionarse realmente; después de o1, los investigadores empezaron a advertir que lo multimodal quizá también necesite cadenas de pensamiento, reflexión, ampliación del espacio de acciones y retroalimentación del entorno para que surja un verdadero momento GPT-4.
\end{importantbox}

\begin{knowledgebox}{Nota sobre la estrategia visual}
Este video es una entrevista con encuadre fijo, sin slides ni demostraciones en pantalla. El texto solo usa la portada para identificar la fuente; todas las imágenes del cuerpo son diagramas conceptuales de elaboración propia, que explican la historia de la transferencia de CV a NLP, la triple separación de la imagen estática, la gap entre generación y comprensión, los defectos de next token, RL/CoT, la Long CoT visual, la colaboración de varios modelos con long context y el aprendizaje en línea.
\end{knowledgebox}

\begin{figure}[H]
\centering
\includegraphics[width=0.96\textwidth]{cv-to-nlp-learning-history.png}
\caption{Historia de cómo la CV aprendió del NLP: de ImageNet/ResNet a ViT, iGPT, MIM y los grandes modelos multimodales. Diagrama conceptual de elaboración propia, basado en la conversación de 00:02:00--00:17:14.}
\end{figure}

\begin{knowledgebox}{Lectura de la figura: el momento GPT de la CV no copió simplemente al NLP}
La figura va de ImageNet y ResNet a ViT/iGPT/MIM y luego a los VLM, y muestra el proceso por el que la CV tomó una y otra vez paradigmas del NLP. La cuestión central es que transferir la arquitectura es fácil, pero transferir las propiedades de los datos y la función objetivo es muy difícil. Los datos de lenguaje provienen por naturaleza de la expresión humana, mientras que la imagen estática no contiene por naturaleza la comprensión humana.
\end{knowledgebox}

\subsection{Ruta de lectura}

Para no quedar abrumados por la terminología, esta sección plantea cuatro preguntas para leer el episodio. Primera: ¿por qué la CV aprendió del NLP durante diez años y aun así no reprodujo GPT con imágenes estáticas? Segunda: ¿por qué la generación y la comprensión de imágenes, aun dentro de un mismo sistema, siguen comportándose como dos modelos? Tercera: ¿por qué, cuanto más grande es el modelo, más fuertes son sus capacidades humanísticas y de conversación, pero su razonamiento matemático puede saltarse pasos y degradarse? Cuarta: ¿por qué la reflexión y los CoT pattern de o1 hicieron que la investigación multimodal volviera a ver un camino?

\noindent\begin{tabularx}{\textwidth}{p{0.19\textwidth}p{0.34\textwidth}X}
\toprule
Pregunta de lectura & Material de la entrevista & Juicio que hay que formarse \\
\midrule
Por qué a la CV le cuesta reproducir GPT & Imagen estática; separación entre generación, comprensión y alineación & La forma de los datos determina el objetivo de aprendizaje; no basta con transferir la arquitectura. \\
Por qué lo multimodal no se ha fusionado & Los modelos de comprensión y de generación mejoran cada uno por su lado, pero no se potencian entre sí & Falta un proceso de pensamiento común y un espacio de acciones controlable. \\
Por qué se degrada el razonamiento & Saltos de pasos en matemáticas de los modelos grandes; defectos de compresión de next token & Comprimir una distribución no equivale a precisión de cálculo. \\
Dónde está el próximo avance & RL, CoT, Long CoT visual, aprendizaje en línea & Lo decisivo está en el espacio de acciones, la retroalimentación y el aprendizaje continuo. \\
\bottomrule
\end{tabularx}

\subsection{Resumen de la sección}

El hilo conductor de EP102 es que lo multimodal no se reduce a «agregar un tipo más de entrada». Para entender el momento GPT-4 de los próximos dos años, hay que comprender a la vez la particularidad de los datos de imagen estática, los defectos de la función objetivo de los modelos de lenguaje, el modo de reflexión que introdujo o1 y el significado del aprendizaje en línea para la próxima generación de Agent.
\section{Diez años de aprendizaje en CV: del scaling de modelos a la migración del paradigma de NLP}

La sección anterior planteó el problema; esta vuelve a la historia. Zhang Xiangyu (张祥雨) sitúa su primera línea de investigación académica en el scaling del aprendizaje profundo posterior a 2012: ImageNet aportó grandes volúmenes de datos, CUDA/GPU aportó capacidad de cómputo, AlexNet demostró que ampliar el modelo mejora notablemente los resultados, y después trabajos como ResNet abrieron el camino a las redes visuales profundas. El núcleo de esta etapa es el model scaling: lograr que los modelos visuales sean más profundos, más grandes y más entrenables.

Más tarde, Transformer, BERT y GPT se convirtieron en el nuevo paradigma de NLP, y la comunidad de CV empezó a trasladar esos métodos. ViT llevó el Transformer a la visión; iGPT intentó modelar imágenes de forma autorregresiva, como si fueran lenguaje; MIM y el aprendizaje contrastivo buscaron aprender representaciones visuales mediante enmascaramiento, aumento de datos e invariancia. Sin embargo, el juicio de Zhang Xiangyu es que muchos de estos métodos solo trasladaron a CV la forma externa de NLP, sin resolver necesariamente la diferencia de fondo de las imágenes estáticas.

\begin{knowledgebox}{Términos clave: qué aprendió CV de NLP}
\noindent\begin{tabularx}{\textwidth}{p{0.18\textwidth}p{0.35\textwidth}X}
\toprule
Término & Mecanismo & Relevancia en este episodio \\
\midrule
ImageNet & Conjunto de datos de imágenes etiquetadas a gran escala & Dio a los modelos visuales, por primera vez, una base de datos estable para el scaling. \\
ResNet & Las conexiones residuales permiten entrenar redes profundas & Demostró que el model scaling visual es viable. \\
ViT & Divide la imagen en patches y la modela con Transformer & La migración de arquitectura tuvo éxito, pero eso no equivale a que la migración de la función objetivo lo tenga. \\
iGPT/MIM & Modelado de imágenes autorregresivo o por enmascaramiento & Intentan replicar el modelado del lenguaje, pero los limitan las propiedades de los datos de imagen. \\
Contrastive Learning & Aprende invariancias a partir de vistas aumentadas & Depende mucho de augmentations diseñadas a mano. \\
\bottomrule
\end{tabularx}
\end{knowledgebox}

\begin{warningbox}{Nota de clase: no confundir «usar Transformer» con el momento GPT}
La advertencia de Zhang Xiangyu es que la arquitectura solo marca la dirección general. ViT, MIM e iGPT le dieron a CV las herramientas de NLP, pero la verdadera pregunta es si con las imágenes estáticas se puede aprender generación, comprensión y alineación con una misma función objetivo.
\end{warningbox}

\subsection{Por qué las imágenes estáticas no son lenguaje natural}

Esta subsección entra en el primer argumento central. El corpus de lenguaje natural lo generan los propios seres humanos, de modo que ya lleva inscritos en la secuencia de tokens los conceptos, las intenciones, la semántica y las relaciones de alineación humanas. Cuando GPT hace modelado generativo sobre el lenguaje, la capacidad de generación, la de comprensión y la alineación con los humanos se refuerzan mutuamente dentro de un mismo objetivo. Las imágenes estáticas son distintas: las producen el mundo natural y el sistema de captura de imagen, y no cambian según cómo las entiendan los seres humanos.

Esto provoca que las tres capacidades queden separadas de forma natural en las imágenes. Generar imágenes permite aprender la distribución de píxeles; comprender imágenes requiere introducir semántica humana; la alineación con los humanos requiere etiquetas, captions, preguntas y respuestas, preferencias o definiciones de tareas. Aunque un modelo genere imágenes realistas, eso no significa que comprenda la escena; aunque un modelo comprenda imágenes, eso no significa que pueda controlar el proceso de generación.

\begin{figure}[H]
\centering
\includegraphics[width=0.92\textwidth]{static-image-three-way-split.png}
\caption{La triple separación de las imágenes estáticas: generación, comprensión y alineación con los humanos están separadas de forma natural en las imágenes estáticas. Diagrama conceptual propio, elaborado a partir de la conversación de 00:18:22--00:24:23.}
\end{figure}

\begin{knowledgebox}{Lectura de la figura: la imagen está ahí; la comprensión la añade el ser humano}
En el centro de la figura está Image, rodeada de generación, comprensión, alineación, naturaleza y tarea. Al leer esta figura hay que retener una idea: la imagen no equivale de forma natural a la semántica humana. Tazas, mesas, acciones y relaciones espaciales necesitan que un marco de tareas humano las inyecte; no se puede obtener una unificación semántica al estilo GPT solo a partir de la distribución de píxeles.
\end{knowledgebox}

\subsection{Resumen de la sección}

Los diez años de avances en CV muestran que los datos, la capacidad de cómputo, la arquitectura y el scaling son importantes; pero también muestran que la visión no es un simple sustituto del lenguaje. Las imágenes estáticas separan generación, comprensión y alineación, y la multimodalidad debe resolver primero esta diferencia en la naturaleza de los datos.
\section{Por qué la generación y la comprensión no se han fusionado realmente}

La sección anterior explicó la diferencia de fondo de las imágenes estáticas; esta sección examina los tropiezos concretos de Zhang Xiangyu (张祥雨) con los sistemas multimodales. Él intentó reunir en un mismo sistema la generación de imágenes, la comprensión de imágenes, la generación de texto y la comprensión de texto: tokenizar las imágenes, procesar lo multimodal con un Transformer de estilo lingüístico y conectar después módulos de generación como diffusion. El resultado fue que el modelo de comprensión se volvió cada vez más capaz y el modelo de generación también, pero al juntarlos no apareció un efecto acumulado en el que 1+1 fuera mayor que 2.

La evidencia más contundente es esta: al quitar la parte de generación, la capacidad de comprensión no se ve afectada; al conectar el modelo de comprensión al modelo de generación, la controlabilidad de la generación sigue siendo deficiente. El modelo puede reconocer que un video contradice el sentido común físico, pero no puede impedirse a sí mismo generar videos que lo contradicen. Esto indica que las dos ramas todavía no comparten un mismo proceso interno de razonamiento.

\begin{figure}[H]
\centering
\includegraphics[width=0.94\textwidth]{generation-understanding-gap.png}
\caption{Generación y comprensión sin fusionar: ambas ramas se fortalecen, pero juntas no dan 1+1 mayor que 2. Diagrama conceptual propio, elaborado a partir de la conversación entre 00:29:10 y 00:38:45.}
\end{figure}

\begin{knowledgebox}{Lectura de la figura: conectar sistemas no equivale a fusionar capacidades}
A la izquierda, el modelo de comprensión destaca en ver imágenes y videos y en alinearlos con el lenguaje; a la derecha, el modelo de generación destaca en producir imágenes y videos. El gap del centro no es una interfaz de ingeniería, sino la imposibilidad de que las capacidades se corrijan mutuamente. Una integración verdadera debería hacer que la comprensión restrinja la generación y que el proceso de generación, a su vez, entrene la comprensión.
\end{knowledgebox}

\subsection{El problema de fondo de la controlabilidad de la generación}

Esta sección desglosa la «baja controlabilidad». Los modelos de generación de imágenes y video pueden aprender patrones visuales de alta probabilidad, pero las escenas complejas exigen respetar restricciones geométricas, físicas, de consistencia de identidad, de continuidad temporal y semánticas. Las indicaciones en lenguaje no siempre logran transmitir con precisión estas restricciones al proceso de generación; tampoco el modelo de comprensión logra necesariamente convertir el «esto no es razonable» en una corrección ejecutable dentro del modelo de generación.

\begin{knowledgebox}{Términos clave: generación, comprensión, alineación, controlabilidad}
\noindent\begin{tabularx}{\textwidth}{p{0.18\textwidth}p{0.36\textwidth}X}
\toprule
Concepto & Significado & Problema en este episodio \\
\midrule
Generación & Producir imágenes, videos o texto & Puede ser realista, pero no necesariamente controlable. \\
Comprensión & Reconocer objetos, relaciones, acciones y semántica & Puede señalar errores, pero no siempre guía la generación. \\
Alineación & Coincidir con la intención y la evaluación humanas & La alineación de imágenes depende del lenguaje, las preferencias y la definición de la tarea. \\
Controlabilidad & Generar de forma estable el resultado buscado bajo restricciones & Las restricciones de múltiples objetos, físicas y temporales son las que más fallan. \\
\bottomrule
\end{tabularx}
\end{knowledgebox}

\begin{warningbox}{El modelo de generación «sabe que está mal», pero aun así lo hace mal}
No es una contradicción superficial. Juzgar si un resultado es erróneo y evitar los errores paso a paso durante la generación son dos capacidades distintas. La primera se parece a calificar; la segunda, a planificar y ejecutar. Lo que les falta a los sistemas multimodales es justamente la cadena de razonamiento que conecte la calificación con la ejecución.
\end{warningbox}

\subsection{Resumen de la sección}

La dificultad de la integración multimodal radica en que la comprensión y la generación no se refuerzan mutuamente con solo conectarlas. Necesitan compartir un proceso intermedio que pueda explicarse, explorarse mediante búsqueda y corregirse; esa es precisamente la razón por la que más adelante entran en la discusión o1, CoT y el Long CoT visual.
\section{Las deficiencias de la Next Token Prediction: comprimir no equivale a precisión de cálculo}

Los problemas anteriores ocurren en la visión; esta sección pasa a los modelos de lenguaje, porque Zhang Xiangyu (张祥雨) considera que en ellos también aparece un cuello de botella similar. Al entrenar modelos grandes, observó un fenómeno extraño: cuanto más grande es el modelo, más fuertes son la conversación general, la inteligencia emocional y la cantidad de conocimiento; pero las capacidades de razonamiento local, como las matemáticas, pueden primero subir, luego estancarse y, si se sigue aumentando la escala, incluso bajar. Un indicio es que los modelos grandes tienden más a «saltarse pasos», en lugar de calcular con cuidado como los modelos pequeños.

Su explicación apunta a la next token prediction. Este objetivo consiste, en esencia, en modelar la distribución de los datos, lo que también puede entenderse como compresión: cuanto mejor predice el modelo el siguiente token, mejor comprime el corpus. Pero comprimir la distribución del corpus humano no equivale a mantener la precisión de cálculo en tareas de matemáticas, lógica, geometría y física. El corpus humano contiene muchos pasos omitidos, intuiciones y elisiones; cuanto mejor comprime un modelo grande, más probable es que aprenda esos atajos.

\[
L_{\mathrm{NTP}}=-\sum_{t}\log p_{\theta}(x_t\mid x_{<t})
\]

Aquí, $x_t$ es el $t$-ésimo token; $x_{<t}$ es el contexto previo; $p_{\theta}$ es la probabilidad que el modelo asigna al siguiente token; cuanto menor es $L_{\mathrm{NTP}}$, mejor se ajusta el modelo a la distribución de entrenamiento. Pero un problema de matemáticas exige que los pasos y el resultado sean correctos, no que la distribución de la salida se parezca a las soluciones de internet.

\begin{figure}[H]
\centering
\includegraphics[width=0.96\textwidth]{next-token-compression-defect.png}
\caption{Las deficiencias de la Next Token Prediction: comprimir la distribución no equivale a precisión de cálculo; cuanto más grande es el modelo, más tiende a saltarse pasos. Diagrama conceptual propio, elaborado a partir de la conversación entre 00:41:11 y 00:50:48.}
\end{figure}

\begin{knowledgebox}{Lectura de la figura: el shortcut es un efecto secundario del aprendizaje de la distribución}
La figura va de Data Distribution a Compression y de ahí a Shortcut, Math y RL. El punto central es el shortcut: si la distribución de entrenamiento contiene muchos pasos omitidos, el modelo aprende una compresión que «se parece a las respuestas humanas», no un procedimiento que «calcula correctamente cada paso».
\end{knowledgebox}

\subsection{Colapso de características: el conflicto entre las características de alta probabilidad y las restricciones precisas}

Esta subsección explica el «colapso de características» mencionado en la entrevista. Los modelos generativos tienden a conservar las características más fáciles de comprimir, más frecuentes y más estables de la distribución de los datos; pero el razonamiento complejo suele depender de unos pocos pasos decisivos. En matemáticas, un paso equivocado invalida todo; en una imagen, si un dedo, la perspectiva o un contacto físico están mal, el conjunto pierde verosimilitud. La similitud de distribución no puede sustituir el cumplimiento de las restricciones.

\begin{figure}[H]
\centering
\includegraphics[width=0.94\textwidth]{feature-collapse.png}
\caption{El fenómeno del colapso de características: los modelos generativos conservan las características de alta probabilidad, pero pierden la precisión de cálculo y de las restricciones. Diagrama conceptual propio, elaborado a partir de la conversación entre 00:45:42 y 00:50:48.}
\end{figure}

\begin{knowledgebox}{Lectura de la figura: parecerse al corpus de entrenamiento no equivale a calcular bien}
El lado izquierdo, «distribución cercana», explica por qué el modelo tiene un lenguaje fluido y un conocimiento amplio; el lado derecho, «precisión insuficiente», explica por qué todavía se equivoca en matemáticas, geometría y física. El razonamiento multimodal necesita restricciones precisas, no solo que la salida se parezca más a las muestras de entrenamiento.
\end{knowledgebox}

\begin{warningbox}{No hay que tomar la tasa de compresión como el límite superior de la inteligencia}
Una mejor next token loss suele mejorar las capacidades generales, pero no es un objetivo suficiente para todas las tareas inteligentes. Las tareas que requieren búsqueda precisa, planificación a largo plazo, verificación externa y retroalimentación de las acciones deben incorporar objetivos y procesos adicionales.
\end{warningbox}

\subsection{Resumen de la sección}

La next token prediction es la base del despegue de los modelos grandes, pero también explica por qué un modelo más grande puede saltarse más pasos. Las matemáticas, la generación visual controlable y el razonamiento multimodal requieren funciones objetivo que vayan más allá de la compresión de la distribución.
\section{Lo que enseña o1: RL, CoT y el pattern de reflexión}

La sección anterior señaló las deficiencias de la next token prediction; esta sección examina cómo entiende o1 Zhang Xiangyu (张祥雨). En la entrevista, su juicio es claro: lo esencial de o1 no es solo que use RL, sino que activa un pattern de cadena de pensamiento. Los nombres de los numerosos algoritmos del RL tradicional no son lo importante; lo importante es que el modelo ya vio durante el preentrenamiento muchos patrones de pensar, reflexionar, verificar y rehacer, y el RL puede reforzarlos.

\begin{figure}[H]
\centering
\includegraphics[width=0.96\textwidth]{rl-cot-pattern.png}
\caption{RL y CoT Pattern: el RL no es solo el nombre de un algoritmo; lo esencial es activar patrones de pensamiento en los que se pueda buscar. Diagrama conceptual propio, elaborado a partir de la conversación de 00:50:48--01:01:52.}
\end{figure}

\begin{knowledgebox}{Lectura de la figura: el RL no es magia; el pattern es la palanca}
En la figura, Pretrain aporta los patrones existentes, el RL los refuerza con retroalimentación sobre el objetivo, CoT despliega los pasos, Reflection corrige la trayectoria y Critical Decision resuelve las bifurcaciones decisivas. Lo que enseña o1 es que el modelo no aprende a pensar de la nada, sino que amplifica los patterns de pensamiento que ya existían, dispersos, en el preentrenamiento.
\end{knowledgebox}

\subsection{Critical Decision: donde un solo token no basta}

Esta subsección explica por qué importa la reflexión. Durante el razonamiento, no todos los tokens son igual de decisivos. Muchos tokens solo dan formato a la expresión; lo que realmente determina el éxito o el fracaso son unas pocas critical decisions: ir a la izquierda o a la derecha, qué fórmula usar, si hace falta volver atrás a verificar, si una conclusión intermedia es confiable. Cuando la complejidad de elegir una rama supera la capacidad expresiva de un solo token, el modelo necesita desplegar más pasos y descomponer el problema de decisión en un espacio en el que se pueda buscar.

\begin{figure}[H]
\centering
\includegraphics[width=0.96\textwidth]{meta-cot-critical-decision.png}
\caption{Meta-CoT y Critical Decision: cuando un solo token no basta para elegir una rama, se necesitan la reflexión y una metacadena de pensamiento. Diagrama conceptual propio, elaborado a partir de la conversación de 01:01:52--01:10:57.}
\end{figure}

\begin{knowledgebox}{Lectura de la figura: Meta-CoT es el CoT del CoT}
De State a Branch, y después a One token?, Reflect y Meta-CoT, la figura expresa «volver a pensar sobre el proceso de pensamiento». El modelo no solo escribe pasos, sino que, en los nodos decisivos, juzga si la trayectoria actual basta, si debe retroceder o si debe cambiar de rama.
\end{knowledgebox}

\begin{importantbox}{Nota de clase: el espacio de acciones importa más que el nombre del algoritmo}
Zhang Xiangyu insiste en que muchas veces lo determinante no es lo misterioso que sea un algoritmo de RL concreto, sino si el espacio de acciones contiene la acción correcta. Sin acciones de reflexión, al modelo le cuesta buscar en los puntos de decisión críticos; con acciones de reflexión, el RL tiene algo que reforzar.
\end{importantbox}

\subsection{Resumen de la sección}

Lo que o1 enseña a la multimodalidad es que la inteligencia no consiste solo en modelos más grandes: también requiere un proceso intermedio en el que se pueda buscar. CoT, la reflexión y Meta-CoT convierten decisiones de un solo token, que no se pueden resolver, en un espacio de acciones que puede descomponerse, explorarse por ensayo y error, y optimizarse.
\section{Long CoT visual: llevar el pensamiento lento al espacio visual}

La sección anterior trató el razonamiento en lenguaje; esta sección vuelve a lo multimodal. Después de estudiar o1, Zhang Xiangyu (张祥雨) volvió a pensar el problema de la poca controlabilidad de la generación visual: muchas tareas de generación visual también superan, en esencia, la complejidad de un razonamiento de un solo paso. Cuando una persona dibuja o comprende relaciones espaciales, no genera una imagen completa a partir de ruido de una sola vez; primero comprende la escena, localiza las partes, revisa las relaciones y después corrige paso a paso. Si lo multimodal solo hace generación de una sola vez, viola con facilidad las restricciones geométricas y físicas.

\begin{figure}[H]
\centering
\includegraphics[width=0.96\textwidth]{visual-long-cot.png}
\caption{Long CoT en el espacio visual: pensar despacio en el espacio visual y trasladar el patrón de reflexión al razonamiento multimodal. Diagrama conceptual propio, elaborado a partir de la conversación de 01:10:57--01:24:07.}
\end{figure}

\begin{knowledgebox}{Lectura de la figura: el CoT visual no consiste en alargar el razonamiento escrito en chino}
La figura va de Image/Video a Visual Actions, Reason, Reflect y Answer. Lo central del Long CoT visual es que el modelo realice acciones en el espacio visual: observar una región, anotar, volver a mirar, comparar, generar borradores auxiliares y corregir relaciones espaciales, en lugar de solo producir explicaciones de texto más largas.
\end{knowledgebox}

\subsection{Por qué la generación también necesita CoT}

Esta sección conecta la comprensión visual con la generación visual. La ruta de Zhang Xiangyu es empezar por el CoT en la comprensión visual, porque en la comprensión es más fácil obtener retroalimentación verificable; pero el objetivo de más largo plazo es que la generación también incorpore una cadena de pensamiento. La generación muy controlable no se dibuja de una sola vez: primero se comprenden las restricciones semánticas y después se aproxima el objetivo mediante generación local, revisión y modificación. En videos complejos esto es todavía más importante, porque la continuidad temporal, la causalidad física y la consistencia de identidad deben mantenerse entre cuadros.

\begin{knowledgebox}{Términos clave: espacio de acciones visuales}
\noindent\begin{tabularx}{\textwidth}{p{0.20\textwidth}p{0.34\textwidth}X}
\toprule
Acción & Significado & Por qué es útil \\
\midrule
Observación local & Mirar solo una parte de la imagen o del video & Reduce la interferencia del contexto y mejora el juicio sobre los detalles. \\
Anotar/trazar recuadros & Generar estructuras intermedias en el espacio visual & Hace explícitas las relaciones espaciales implícitas. \\
Volver a mirar/comparar & Contrastar distintas regiones o momentos & Detecta inconsistencias y errores físicos. \\
Generar borradores & Generar primero resultados intermedios que se puedan revisar & Permite que el proceso de generación reciba retroalimentación y se corrija. \\
Reflexión y corrección & Rehacer el camino tras detectar un error & Convierte la tarea visual en un proceso de búsqueda. \\
\bottomrule
\end{tabularx}
\end{knowledgebox}

\subsection{Dos frentes: corpus de preentrenamiento y espacio de acciones}

Lo anterior mostró que el Long CoT visual necesita acciones; esta sección concentra la ruta en dos frentes. La ruta futura de lo multimodal no es una sola línea. El primer frente es ampliar el corpus de preentrenamiento: más datos de alta calidad de imagen y texto, de video, de procesos de generación, de procesos de interacción y de razonamiento visual. El segundo frente es ampliar el espacio de acciones: que el modelo no solo pueda ver y hablar, sino también observar regiones, dibujar, generar, revisar, llamar herramientas, llamar submodelos y recibir retroalimentación. Solo si los datos y el espacio de acciones se amplían juntos, el pensamiento visual tendrá la oportunidad de acercarse al salto que dio el razonamiento en lenguaje.

\begin{figure}[H]
\centering
\includegraphics[width=0.94\textwidth]{two-leg-roadmap.png}
\caption{Los dos frentes futuros de lo multimodal: ampliar el corpus de preentrenamiento y, al mismo tiempo, ampliar el espacio de acciones. Diagrama conceptual propio, elaborado a partir de la conversación de 01:31:31--01:45:42.}
\end{figure}

\begin{knowledgebox}{Lectura de la figura: los datos resuelven «haberlo visto»; el espacio de acciones resuelve «saber hacerlo»}
A la izquierda, Pretraining Data permite que el modelo vea más procesos visuales y de generación; a la derecha, Action Space le da al modelo más acciones ejecutables dentro de la tarea. Es muy probable que el «momento GPT-4» de lo multimodal requiera que ambas cosas ocurran a la vez, en lugar de seguir solo acumulando pares de imagen y texto.
\end{knowledgebox}

\subsection{Línea de tiempo de la ruta de investigación}

Esta sección condensa los conceptos anteriores en una línea de tiempo, para que el lector entienda este episodio desde la historia de la investigación y no desde las palabras de moda. El relato de Zhang Xiangyu pasa aproximadamente por cuatro etapas: la primera es el CV scaling, hacer más profundos y más grandes modelos como ResNet; la segunda es la CV aprendiendo de NLP, con el intento de trasladar a las imágenes el Transformer, BERT/GPT, MIM y las ideas autorregresivas; la tercera es el tropiezo de la integración multimodal, al descubrir que la generación y la comprensión no se fortalecían de verdad entre sí; la cuarta es volver a entender el razonamiento después de o1 y llevar de regreso a la visión el RL, el CoT, la reflexión y el espacio de acciones.

\noindent\begin{tabularx}{\textwidth}{p{0.16\textwidth}p{0.31\textwidth}X}
\toprule
Etapa & Acción central & Lección obtenida \\
\midrule
2012--2016 & ImageNet, AlexNet y ResNet demuestran el model scaling en visión & Ampliar a la vez datos, cómputo y modelo abre el reconocimiento visual. \\
2018--2022 & ViT, iGPT, MIM, aprendizaje contrastivo y otros aprenden de NLP & La transferencia de arquitecturas funciona, pero la función objetivo y las propiedades de los datos no se resolvieron de fondo. \\
2023--2024 & Entrenar grandes modelos multimodales e intentar integrar generación y comprensión & La comprensión y la generación se fortalecen cada una por su lado, pero no se fusionan automáticamente. \\
Después de o1 & Volver a ver lo multimodal desde el RL, el CoT, la reflexión y el espacio de acciones & La visión también necesita pensamiento lento y procesos intermedios en los que se pueda buscar. \\
Próximos dos años & Long context, aprendizaje en línea, aprendizaje autónomo y modelos del mundo & El momento GPT-4 podría venir de la retroalimentación del entorno y de la colaboración entre sistemas. \\
\bottomrule
\end{tabularx}

\begin{importantbox}{El docente enfatiza: la desorientación misma es evidencia sobre la ruta}
En la entrevista aparecen una y otra vez expresiones como «pasé más de medio año muy desorientado» o «cuando salió o1 lo vi claro». Esa desorientación no es ruido, sino evidencia sobre la ruta técnica: si una dirección sigue sin lograr 1+1>2 aunque se agreguen datos, se agrande el modelo o se conecten módulos, lo que falta es un nuevo objetivo de entrenamiento o un nuevo espacio de acciones.
\end{importantbox}

\subsection{Resumen de la sección}

El Long CoT visual es la idea multimodal nueva más importante de este episodio: trasladar a la visión las ideas de reflexión y de espacio de acciones de o1, para que el modelo piense despacio en el espacio visual. No se trata de alargar la cadena de texto, sino de darle al modelo multimodal más acciones visuales verificables, en las que se pueda buscar y que se puedan corregir.
\section{Long Context, colaboración entre múltiples modelos y aprendizaje en línea}

Las secciones anteriores trataron la comprensión y la generación multimodales; esta sección aborda los dos momentos GPT-4 de los próximos dos años. Zhang Xiangyu (张祥雨) considera que el long context es importante, pero que introducir hoy un contexto extremadamente largo directamente en un solo Transformer plantea problemas: la atención se dispersa; cuanto más largo es el contexto, más información similar contiene y más fácil es que el modelo se distraiga. La inteligencia verdadera no consiste en recordarlo todo sin pérdida, sino en poder comprimir, olvidar, volver a consultar el libro y desplazar la atención.

\begin{figure}[H]
\centering
\includegraphics[width=0.96\textwidth]{long-context-multi-model.png}
\caption{Long Context y colaboración entre múltiples modelos: no hace falta introducir el libro entero en un solo contexto, sino combinar una planificación global con una ejecución local. Diagrama conceptual propio, elaborado a partir de la conversación de 01:46:56--02:07:09.}
\end{figure}

\begin{knowledgebox}{Lectura de la figura: la planificación global y la ejecución local pueden separarse}
En la figura, Global Planner conserva la impresión de conjunto, Retrieve/Turn page se encarga de volver a consultar el libro, Local Worker realiza el razonamiento local, Feedback devuelve los resultados y Context Shift cambia de escenario. Esto se acerca más al mecanismo de atención humano que introducir todos los tokens en un mismo contexto.
\end{knowledgebox}

\subsection{Por qué la arquitectura no es lo esencial}

Esta subsección explica la idea de que «Linear Transformer no es lo esencial». Zhang Xiangyu no niega el valor de la atención lineal ni de las arquitecturas RNN-like; lo que sostiene es que, si el objetivo es solo soportar contextos extremadamente largos a fuerza de arquitectura, quizá no se llega al núcleo del problema. El núcleo del contexto largo es cómo comprimir, cómo recuperar información, cómo cambiar de escenario y cómo hacer que varios modelos o módulos colaboren, no solo reducir la complejidad de la atención por debajo de la cuadrática.

\begin{warningbox}{No confundir la búsqueda de una aguja en un pajar con la inteligencia misma}
Para superar las pruebas de retrieve, el modelo puede aprender el bias de «no se puede olvidar nada». Pero cuando una persona escribe un resumen después de leer un libro, recuerda la estructura y luego vuelve al libro para encontrar las frases que necesita. Un sistema de long context realmente útil debería permitir olvidar, resumir, pasar páginas, revisar localmente y cambiar de contexto.
\end{warningbox}

\begin{knowledgebox}{Términos clave: los dos tipos de problemas del long context}
\noindent\begin{tabularx}{\textwidth}{p{0.20\textwidth}p{0.35\textwidth}X}
\toprule
Problema & Significado & Valoración en este episodio \\
\midrule
Complejidad computacional & El costo que genera la atención al crecer con la longitud & Importante, pero no es el único problema. \\
Modelado de la inteligencia & Cómo comprimir, olvidar, desplazar la atención y revisar localmente & Más cercano a la verdadera inteligencia de largo alcance. \\
Colaboración entre múltiples modelos & Repartir el procesamiento del contexto entre planner y worker & Puede reducir la interferencia y hacer que el sistema sea entrenable. \\
Aislamiento de escenarios & Al resolver el segundo problema no hace falta cargar con todo el contexto del primero & Evita que el contexto largo contamine el razonamiento actual. \\
\bottomrule
\end{tabularx}
\end{knowledgebox}

\subsection{Aprendizaje en línea y aprendizaje autónomo}

Si el long context es un posible momento GPT-4, el otro es el aprendizaje en línea o aprendizaje autónomo. Zhang Xiangyu lo relaciona con los Agent: hoy muchos Agent son orquestación de herramientas en la capa de aplicación, pero un Agent más fundamental debería poder actuar de forma continua en un entorno, recibir retroalimentación y mejorar su estrategia. Aquí el scaling ya no es solo model scaling o data scaling, sino environmental scaling: si se puede proporcionar retroalimentación del entorno suficientemente abundante, realista y verificable.

\begin{figure}[H]
\centering
\includegraphics[width=0.96\textwidth]{online-learning-agent.png}
\caption{Aprendizaje en línea y aprendizaje autónomo: el próximo momento GPT-4 podría surgir de la retroalimentación del entorno y de la automejora continua. Diagrama conceptual propio, elaborado a partir de la conversación de 02:08:30--02:25:00.}
\end{figure}

\begin{knowledgebox}{Lectura de la figura: la retroalimentación del entorno es el nuevo objeto del scaling}
En la figura, Environment, Agent, Feedback, RL/IRL y Online Learning forman un ciclo cerrado. El modelo ya no aprende solo de corpus fuera de línea, sino que se actualiza continuamente con tareas, herramientas, el mundo y la retroalimentación humana. Las dificultades son la escala del entorno, la calidad de la retroalimentación y los límites de seguridad.
\end{knowledgebox}

\begin{importantbox}{Nota de clase: una aplicación de Agent y un Agent de aprendizaje autónomo no son lo mismo}
Hoy muchos productos de Agent consisten en orquestación de herramientas y automatización de flujos de trabajo; el Agent más fundamental que discute Zhang Xiangyu es un sistema de aprendizaje capaz de mejorar continuamente a partir de la retroalimentación del entorno. El primero es una forma de aplicación; el segundo, un paradigma de entrenamiento y un mecanismo de inteligencia.
\end{importantbox}

\subsection{Modelos del mundo y la robótica que se adelanta}

La entrevista termina vinculando la multimodalidad con la robótica. Las personas no tienen un órgano de generación visual, pero sí un modelo del mundo: podemos predecir mentalmente el espacio, la física y las consecuencias de las acciones. Sigue siendo una pregunta abierta si un sistema de aprendizaje automático debe obtener un modelo del mundo mediante entrenamiento generativo; pero si un robot ha de actuar en el mundo físico, no puede prescindir de la comprensión visual, la retroalimentación de las acciones ni el modelo del mundo.

\begin{figure}[H]
\centering
\includegraphics[width=0.94\textwidth]{world-model-robotics-gap.png}
\caption{Modelos del mundo y la robótica que se adelanta: las personas no tienen órganos de generación, pero sí un modelo del mundo; los robots no pueden saltarse la comprensión visual. Diagrama conceptual propio, elaborado a partir de la conversación de 02:25:00--02:28:44.}
\end{figure}

\begin{knowledgebox}{Lectura de la figura: la robótica no puede apostar solo por el cuerpo físico}
A la izquierda, World Model se encarga del razonamiento espacial, el sentido común físico y la predicción del futuro; a la derecha, Robotics requiere que la visión, la acción, la retroalimentación y el cuerpo del robot superen juntos el umbral. Si la industria robótica se salta el razonamiento visual y el modelo del mundo y depende solo del cuerpo del robot y de la teleoperación, es muy probable que se adelante a lo que la inteligencia disponible permite.
\end{knowledgebox}

\subsection{Resumen de la sección}

Las dos direcciones clave de avance para el futuro son el modelado de una inteligencia real en el long context y el ciclo cerrado de retroalimentación del entorno en el aprendizaje en línea o autónomo. La primera exige que el modelo aprenda a comprimir, volver a consultar el libro y cambiar de escenario; la segunda, que el modelo mejore continuamente mientras actúa.
\section{Síntesis y ampliación}

Esta sección condensa el EP102 en siete conclusiones. Primera: el momento GPT de la CV no puede copiar sin más el de la NLP, porque las imágenes estáticas no traen de forma natural semántica humana ni alineación. Segunda: la dificultad central de la integración multimodal es que la generación y la comprensión no pueden corregirse entre sí. Tercera: la next token prediction es la base del despegue de los grandes modelos y también una de las causas de los saltos de pasos en el razonamiento y de la falta de precisión. Cuarta: la aportación principal de o1 es usar RL para activar la CoT, la reflexión y la búsqueda de critical decisions. Quinta: la Long CoT visual traslada el pensamiento lento al espacio visual. Sexta: el long context no es solo un problema de complejidad de la arquitectura, sino sobre todo de compresión, recuperación, aislamiento de escenarios y colaboración entre varios modelos. Séptima: el siguiente momento Agent, más profundo, podría venir del aprendizaje en línea y del environmental scaling.

\begin{importantbox}{El EP102 dentro de la serie de IA e internet de Zhang Xiaojun (张小珺)}
El EP102 es una de las entrevistas sobre la ruta multimodal con mayor contenido técnico de este lote. Se conecta con el informe técnico sobre Agents del EP110, con el agentic LLM del EP113 y con la filosofía de investigación sobre Agents del EP115; además, ofrece para los problemas de inteligencia corporizada de los EP106 y EP109 una explicación previa sobre razonamiento visual y modelos del mundo.
\end{importantbox}

\subsection{Glosario general}

Las secciones anteriores ya explicaron por separado VLM, CoT, espacio de acciones, aprendizaje en línea y modelo del mundo. Esta sección reúne esos términos para que el lector pueda repasar la cadena argumentativa de todo el episodio: desde las propiedades de los datos de las imágenes estáticas, pasando por la función objetivo de la next token prediction, hasta la reflexión al estilo o1, la Long CoT visual y el aprendizaje en línea.

\noindent\begin{tabularx}{\textwidth}{p{0.18\textwidth}p{0.36\textwidth}X}
\toprule
Término & Definición breve & Papel en este episodio \\
\midrule
VLM & Vision-Language Model, modelo de visión y lenguaje & Forma principal de modelo que conecta imágenes, video y texto. \\
MIM & Masked Image Modeling, modelado de imágenes enmascaradas & Ruta representativa de la CV hacia el aprendizaje por preentrenamiento al estilo BERT. \\
Next Token Prediction & Objetivo de entrenamiento que predice el siguiente token & Explica el origen común de la potencia de los modelos de lenguaje y de su defecto de saltarse pasos. \\
CoT & Chain of Thought, cadena de pensamiento & Despliega una decisión de un solo paso en un proceso en el que se puede buscar. \\
Meta-CoT & Repensar y reflexionar sobre la CoT & Atiende las critical decisions y la elección de rutas. \\
Action Space & Conjunto de acciones que el modelo puede elegir en una tarea & Determina si el RL puede encontrar mediante búsqueda la conducta correcta. \\
Online Learning & Aprendizaje en línea: el sistema se actualiza de forma continua mientras se usa & Podría ser la clave de la siguiente generación de Agents. \\
World Model & Modelo interno que predice cómo cambia el estado del mundo & Conecta lo multimodal, la robótica y la inteligencia corporizada. \\
\bottomrule
\end{tabularx}

\subsection{Preguntas para observar en adelante}

Esta sección no enumera preguntas genéricas: convierte los juicios de la entrevista sobre las rutas técnicas en puntos de observación que se pueden seguir. En los próximos dos años, si de verdad aparece un nuevo momento GPT-4 en lo multimodal, las señales deberían verse en estos puntos: no solo que los modelos suban en los benchmarks, sino que surjan nuevas capacidades de sistema en razonamiento visual, controlabilidad de la generación, long context y aprendizaje en línea.

\begin{enumerate}
\item ¿Puede la Long CoT visual mejorar de forma notable el razonamiento multimodal en benchmarks reales, y no solo generar explicaciones más largas?
\item ¿Dará un salto cualitativo la controlabilidad de los modelos generativos gracias a la CoT visual, la edición local y el entrenamiento con retroalimentación?
\item ¿El long context seguirá dependiendo de ampliar la ventana de un solo modelo, o se moverá hacia esquemas planner/worker, la colaboración entre varios modelos y la consulta mediante herramientas, como quien hojea un libro?
\item ¿Cómo resolverá el aprendizaje en línea los problemas de seguridad, contaminación de datos, retroalimentación escasa y costo de la actualización continua?
\item ¿Podrá la robótica, antes de adelantarse en el hardware y en el cuerpo del robot, completar las capacidades básicas de razonamiento visual, modelo del mundo y retroalimentación de acciones?
\end{enumerate}

\subsection{Lecturas adicionales}

\begin{itemize}
\item Si le interesan el entrenamiento de Agents y la ingeniería de contexto, puede compararlo con las notas del EP110 sobre los informes técnicos de Kimi K2 / ChatGPT Agent / Qwen3-Coder.
\item Si le interesan el Agentic LLM y el escalamiento en tiempo de prueba, puede compararlo con las notas del EP113, la entrevista a Yang Zhilin (杨植麟) sobre Kimi K2.
\item Si le interesan los problemas de datos y de modelos del mundo en la inteligencia corporizada, puede compararlo con las entrevistas del EP106 a Wang He (王鹤), del EP109 a Xie Chen (谢晨) y del EP121 a Tan Jie (谭捷).
\end{itemize}

\end{document}
